
\documentclass[11pt,a4paper]{article}

\usepackage[T1]{fontenc}
\usepackage{lmodern}
\usepackage{microtype}
\usepackage[margin=28mm]{geometry}
\usepackage{amsmath,amssymb,amsthm,mathtools}
\usepackage{enumitem}
\usepackage[hidelinks]{hyperref}
\usepackage{fancyhdr}

\setlength{\parindent}{0pt}
\setlength{\parskip}{0.65em}
\setlist[enumerate]{leftmargin=2em,itemsep=0.2em,topsep=0.2em}
\setlist[itemize]{leftmargin=2em,itemsep=0.2em,topsep=0.2em}
\allowdisplaybreaks[2]
\pagestyle{fancy}
\fancyhf{}
\fancyhead[L]{\small A Note on Traces}
\fancyhead[R]{\small 10 November 2026}
\fancyfoot[C]{\thepage}
\renewcommand{\headrulewidth}{0.3pt}

% Colour Section Titles
\usepackage[dvipsnames]{xcolor}
\usepackage{titlesec}
\usepackage{enumitem}

% Current heading colour (stored as a name)
\def\currentheadingcolor{Plum}

% Section heading: Plum
\titleformat{\section}
  {\normalfont\Large\bfseries\color{Plum}}
  {\thesection}{1em}{}

% Subsection heading: Orchid
\titleformat{\subsection}
  {\normalfont\large\bfseries\color{Orchid}}
  {\thesubsection}{1em}{}

% For highlightng text
\newcommand{\skyblue}[1]{\textcolor{SkyBlue}{#1}}

\newcommand{\Om}{\Omega}
\newcommand{\Gam}{\Gamma}
\newcommand{\dd}{\,\mathrm{d}}
\newcommand{\tr}{\operatorname{tr}}
\newcommand{\Div}{\operatorname{div}}
\newcommand{\Curl}{\operatorname{curl}}
\newcommand{\grad}{\nabla}
\newcommand{\Hdiv}{H(\Div;\Om)}
\newcommand{\Hcurl}{H(\Curl;\Om)}
\newcommand{\HDelta}{H(\Delta;\Om)}
\newcommand{\dual}[2]{\left\langle #1,#2\right\rangle}

\title{\textbf{A Note on Traces\\\large Weak Boundary Traces via Green's Identities}}
\author{}
\date{}

\begin{document}
\maketitle
\vspace{-1.5em}

\section{Disclaimer}
These are notes written by chat (edited by me) based off of my handwritten notes. The relevant references are \cite{BoffiTrace} and \cite{BeeckTrace}.

\section{Motivation}

Functions in $H^1(\Om)$ are not necessarily continuous. Thus, it is not
immediately clear what it means to evaluate $u\in H^1(\Om)$ on $\partial\Om$.

Another subtlety is that $H^1$ functions are equivalence classes of functions
that agree almost everywhere in $\Om$. However, the boundary of a Lipschitz
domain has zero $d$-dimensional Lebesgue measure. We could therefore change
the values of a representative on the boundary without changing its
equivalence class.

The solution is to introduce a \emph{trace operator}, which gives a
well-defined meaning to boundary values in suitable boundary function spaces.

\section{The trace of \texorpdfstring{$H^1(\Om)$}{H1} functions}

When $\partial\Om$ is sufficiently regular (for example, when $\Om$ is a
bounded Lipschitz domain), the classical restriction of smooth functions to
the boundary is re-thought of using a continuous trace operator
\[
    \tr_0:H^1(\Om)\longrightarrow H^{1/2}(\partial\Om),
    \qquad \tr_0 u = u|_{\partial\Om}
    \quad\text{for smooth }u.
\]
The traces of functions in $H^1(\Om)$ form $H^{1/2}(\partial\Om)$, which
is a proper dense subspace of $L^2(\partial\Om)$.
In particular, $\tr_0$ is surjective and possesses a continuous lifting. (what does that mean?)
% there exists a bounded linear operator
% \[
%     E:H^{1/2}(\partial\Om)\longrightarrow H^1(\Om)
%     \qquad\text{such that}\qquad \tr_0(Eg)=g.
% \]
% We will use this lifting when showing that weak boundary traces are bounded.

\section{The \texorpdfstring{$H(\Div;\Om)$}{H(div)} space}

Recall that
\[
    \Hdiv
    =\bigl\{\mathbf q\in L^2(\Om)^d:
                   \Div\mathbf q\in L^2(\Om)\bigr\},
\]
with norm
\[
    \|\mathbf q\|_{H(\Div;\Om)}^2
    =\|\mathbf q\|_{L^2(\Om)}^2
     +\|\Div\mathbf q\|_{L^2(\Om)}^2.
\]

For $\mathbf q\in\Hdiv$, we define its \emph{normal trace} through Green's
identity. In particular, for sufficiently smooth fields,
\[
    \int_{\partial\Om}(\mathbf q\cdot\boldsymbol\nu)v\dd s
    =\int_\Om(\Div\mathbf q)v\dd x
     +\int_\Om\mathbf q\cdot\grad v\dd x,
\]
where $\boldsymbol\nu$ is the outward unit normal. This motivates the
variational definition
\begin{equation}\label{eq:hdiv-trace}
    \dual{\gamma_n\mathbf q}{\tr_0v}_{H^{-1/2},H^{1/2}}
    :=\int_\Om(\Div\mathbf q)v\dd x
      +\int_\Om\mathbf q\cdot\grad v\dd x,
    \qquad v\in H^1(\Om).
\end{equation}
This expression is used to define the normal trace:
\[
    \gamma_n:\Hdiv\longrightarrow H^{-1/2}(\partial\Om),
    \qquad \gamma_n\mathbf q=\mathbf q\cdot\boldsymbol\nu
    \quad\text{for smooth }\mathbf q.
\]
which is well-defined, bounded, and surjective when $\Omega$ has a Lipschitz continuous boundary. (any conditions I have missed?)

% The expression on the right depends only on $\tr_0v$: if $v\in H^1_0(\Om)$,
% it vanishes by the definition of the distributional divergence and density
% of $C_c^\infty(\Om)$ in $H^1_0(\Om)$. Together with the trace lifting and
% Cauchy--Schwarz, this gives a well-defined bounded operator
% \[
%     \gamma_n:\Hdiv\longrightarrow H^{-1/2}(\partial\Om),
%     \qquad \gamma_n\mathbf q=\mathbf q\cdot\boldsymbol\nu
%     \quad\text{for smooth }\mathbf q.
% \]
% For bounded Lipschitz domains, this operator is also surjective.

\section{The \texorpdfstring{$H(\Curl;\Om)$}{H(curl)} space}

In three dimensions, recall that
\[
    \Hcurl
    =\bigl\{\mathbf x\in (L^2(\Om))^3:
                   \Curl\mathbf x\in (L^2(\Om))^3\bigr\},
\]
with norm
\[
    \|\mathbf x\|_{H(\Curl;\Om)}^2
     =\|\mathbf x\|_{L^2(\Om)}^2
      +\|\Curl\mathbf x\|_{L^2(\Om)}^2.
\]

To define the tangential trace, we start from integration by parts:
\begin{equation}\label{eq:curl-green}
    \int_\Om \mathbf x\cdot\Curl\boldsymbol\Phi\dd x
    -\int_\Om (\Curl\mathbf x)\cdot\boldsymbol\Phi\dd x
    =\int_{\partial\Om}
        (\mathbf x\times\boldsymbol\nu)\cdot\boldsymbol\Phi\dd s.
\end{equation}
The terms on the left-hand-side are meaningful for $\mathbf x\in\Hcurl$ and
$\boldsymbol\Phi\in H^1(\Om)^3$. We therefore use them to define the
\emph{tangential trace} $\gamma_t\mathbf x$, by setting
\begin{equation}\label{eq:hcurl-trace}
    \dual{\gamma_t\mathbf x}{\tr_0\boldsymbol\Phi}
    :=\int_\Om\mathbf x\cdot\Curl\boldsymbol\Phi\dd x
      -\int_\Om(\Curl\mathbf x)\cdot\boldsymbol\Phi\dd x.
\end{equation}
As before, the right-hand side depends only on the boundary trace of the
test function. Consequently,
\[
    \gamma_t:\Hcurl\longrightarrow (H^{-1/2}(\partial\Om))^3,
    \qquad \gamma_t\mathbf x=(\mathbf x\times\boldsymbol\nu)|_{\partial \Om}
       \quad\text{for smooth }\mathbf x,
\]
is well-defined and continuous.

Surjectivity is more complicated. The tangential trace of an $H(\Curl; \Om)$
field belongs to a more specific space than $(H^{-1/2}(\partial\Om))^3$.
To see this, notice that
\[
    (\gamma_t\mathbf x)\cdot\boldsymbol\nu=0.
\]
Thus, the range of $\gamma_t$ lies in the tangent plane. However, $(H^{-1/2}(\partial\Om))^3$ has elements with nonzero normal components. 


In particular, for smooth domains, the range of $\gamma_t$ is the following space:
\[
\begin{split}
    H^{-1/2}(\Div_{\partial \Om};\partial \Om)
    :=\bigl\{\mathbf q\in H^{-1/2}(\partial \Om)^3:\;&
      \mathbf q\cdot\boldsymbol\nu=0,\\[-0.2em]
     &\Div_{\partial \Om} \mathbf q\in H^{-1/2}(\partial \Om)\bigr\},
\end{split}
\]
where $\Div_{\partial \Om}$ denotes \emph{surface divergence}. For example, on the
flat surface $\partial \Om =\{(x_1,x_2,x_3):x_3=0\}$, a tangential field
$\mathbf q=(q_1,q_2,0)$ has
\[
    \Div_{\partial \Om} \mathbf q
     =\frac{\partial q_1}{\partial x_1}
      +\frac{\partial q_2}{\partial x_2}.
\]
% With this range space, the tangential trace is surjective (for suitably
% regular $\Gamma$).

Before moving on, we note that care is required when the domain is polyhedral rather than smooth.
On polyhedral domains, the normal jumps across edges. As a result, $q\cdot\boldsymbol\nu$ is not well-defined for $\mathbf q\in H^{-1/2}(\partial \Om)^3$ and our definition of $H^{-1/2}(\Div_{\partial \Om};\partial \Om)$ breaks down.

% so the preceding simple description using
% multiplication by a smooth normal field is no longer appropriate; the
% correct trace spaces must account for the geometry of the boundary.
% The tangential trace operator itself nevertheless remains well-defined.

\section{Example: the Laplacian}

Consider
\[
    -\Delta u=f\quad\text{in }\Om,
    \qquad u\in H^1(\Om),\quad f\in L^2(\Om).
\]
For sufficiently smooth functions, integration by parts yields
\begin{equation}\label{eq:poisson-green}
    \int_\Om fv\dd x
      =\int_\Om\grad u\cdot\grad v\dd x
       -\int_{\partial\Om}(\grad u\cdot\boldsymbol\nu)v\dd s,
    \qquad v\in H^1(\Om).
\end{equation}
Notice that since $\grad u\in L^2(\Om)^d$, its classical boundary restriction need not
exist. However, the PDE supplies additional information we need to define a trace. In particular, since
\[
    \Div(\grad u)=\Delta u=-f\in L^2(\Om),
\]
we know that $\grad u\in\Hdiv.$ As a result,  we can use the normal trace from
\eqref{eq:hdiv-trace} to define
\begin{equation}\label{eq:poisson-normal}
    \dual{\partial_\nu u}{\tr_0v}
       :=\int_\Om\grad u\cdot\grad v\dd x
         -\int_\Om fv\dd x
       =\int_\Om\grad u\cdot\grad v\dd x
         +\int_\Om(\Delta u)v\dd x,
    \qquad v\in H^1(\Om).
\end{equation}
Here $\partial_\nu u=\gamma_n(\grad u)$.

\subsection*{Well-defined, Uniqueness, and Bounded}

We first show that the right-hand side of \eqref{eq:poisson-normal} depends
only on $\tr_0v$. Suppose $v_1,v_2\in H^1(\Om)$ satisfy
$\tr_0v_1=\tr_0v_2$. Then $w=v_1-v_2\in H^1_0(\Om)$.

Since $H^1_0(\Om)$ is the closure of $C_c^\infty(\Om)$ in the $H^1(\Om)$ norm, the
following identity (for the distributional Laplacian), 
\[
    \int_\Om(\Delta u)\varphi\dd x
        =-\int_\Om\grad u\cdot\grad\varphi\dd x,
    \qquad \varphi\in C_c^\infty(\Om),
\]
extends to every $w\in H^1_0(\Om)$. Therefore,
\[
    \int_\Om\grad u\cdot\grad w\dd x
       +\int_\Om(\Delta u)w\dd x=0.
\]
It follows that the two choices $v_1$ and $v_2$ give the same value in
\eqref{eq:poisson-normal}.

Moreover, $\tr_0:H^1(\Om)\to H^{1/2}(\partial\Om)$ is surjective. Hence,
for every $g\in H^{1/2}(\partial\Om)$, we may choose $v\in H^1(\Om)$
with $\tr_0v=g$ and set
\[
    \dual{\partial_\nu u}{g}
      =\int_\Om\grad u\cdot\grad v\dd x
        +\int_\Om(\Delta u)v\dd x.
\]
As a result, the trace operator $\partial_\nu u$ is well-defined and unique. Furthermore, $\partial_\nu u$ is bounded, making it an element of $H^{-1/2}(\partial\Om).$

% \subsection*{Boundedness}

% By Cauchy--Schwarz,
% \[
%     |\dual{\partial_\nu u}{\tr_0v}|
%       \leq \bigl(\|\grad u\|_{L^2(\Om)}
%                   +\|\Delta u\|_{L^2(\Om)}\bigr)
%               \|v\|_{H^1(\Om)}.
% \]
% To obtain a bound in the boundary norm, take $v=Eg$, where $E$ is the
% continuous lifting of the trace. Then
% \[
%     |\dual{\partial_\nu u}{g}|
%       \leq C\bigl(\|\grad u\|_{L^2(\Om)}
%                  +\|\Delta u\|_{L^2(\Om)}\bigr)
%               \|g\|_{H^{1/2}(\partial\Om)}.
% \]
% Consequently,
% \[
%     \partial_\nu u\in H^{-1/2}(\partial\Om).
% \]

\newpage

\section{A biharmonic-type problem}

We now apply the same idea to a fourth-order problem, for which there are
two boundary traces to define. The starting point is Green's second identity.

\subsection{Model problem and boundary quantities}

Let $\Om\subset\mathbb R^d$, with $d\in\{2,3\}$, be a bounded convex
domain with smooth boundary, and let $\boldsymbol\nu$ denote its outward
unit normal. Consider the model equation
\begin{equation}\label{eq:biharmonic-model}
    \alpha\Delta^2u
      +\beta\,\Div\!\left(\grad\bigl(\mathbf n^{\mathsf T}
                   (\nabla^2u)\mathbf n\bigr)\right)
      -\Delta u-\kappa^2u=f
    \qquad\text{in }\Om,
\end{equation}
with the (unspecified) boundary condition $\mathcal Bu=0$ on $\partial\Om$.
Here $f\in L^2(\Om)$, $\alpha,\beta\geq0$ are constant coefficients,
$\nabla^2u$ denotes the Hessian of $u$, and $\mathbf n$ is a piecewise
constant unit vector field describing the local average orientation of
molecules in the fluid. In particular, $\mathbf n$ is distinct from the
outward boundary normal $\boldsymbol\nu$.

We are interested in the boundary quantities
\begin{align}
    T_0u&:=\alpha\Delta u
                 +\beta\,\mathbf n^{\mathsf T}(\nabla^2u)\mathbf n,
                    \label{eq:T0formal}\\
    T_1u&:=\alpha\grad(\Delta u)\cdot\boldsymbol\nu
       +\beta\grad\!\left(\mathbf n^{\mathsf T}
                     (\nabla^2u)\mathbf n\right)\cdot\boldsymbol\nu
       =\partial_\nu(T_0u).
                    \label{eq:T1formal}
\end{align}
% where the expressions in \eqref{eq:T1formal} are initially understood
% formally for sufficiently smooth functions.

As in the Poisson example, the PDE yields additional regularity. Assuming
$u\in H^2(\Om),$ \eqref{eq:biharmonic-model} gives
\[
    \Delta(T_0u)
      =\alpha\Delta^2u
       +\beta\Delta\!\left(\mathbf n^{\mathsf T}
                          (\nabla^2u)\mathbf n\right)
      =f+\Delta u+\kappa^2u\in L^2(\Om).
\]
Since $u\in H^2(\Om)$ also implies $T_0u\in L^2(\Om)$, we find that
\[
    T_0u\in\HDelta
       :=\bigl\{q\in L^2(\Om):\Delta q\in L^2(\Om)\bigr\}.
\]
This is just enough regularity to employ Green's second identity and to thus define the needed boundary traces. In particular, for sufficiently smooth $q$ and $w$, Green's second identity is as follows:
\begin{equation}\label{eq:greens-second}
    \int_\Om(\Delta q)w\dd x
      =\int_\Om q\,\Delta w\dd x
       +\int_{\partial\Om}(\partial_\nu q)w\dd s
       -\int_{\partial\Om}q\,\partial_\nu w\dd s.
\end{equation}
Taking $q=T_0u$, the two boundary integrals contain $T_1u=\partial_\nu(T_0u)$
and $T_0u$. As before, we use the well-define volume integrals in this identity to define the not-so-well-defined
boundary traces. 

\subsection{The trace \texorpdfstring{$T_0u\in H^{-1/2}(\partial\Om)$}{T0}}

Starting from \eqref{eq:greens-second}, take $w\in H^2(\Om)$ with
$\tr_0w=0$. This removes the boundary term involving $T_1u$ and motivates
\begin{equation}\label{eq:T0weak}
    \dual{T_0u}{\tr_\nu w}_{H^{-1/2},H^{1/2}}
       :=\int_\Om(T_0u)\Delta w\dd x
         -\int_\Om\bigl(\Delta(T_0u)\bigr)w\dd x,
    % \quad
    % \begin{gathered}
    %    w\in H^2(\Om),\\[-0.2em]
    %    \tr_0w=0,
    % \end{gathered}
\end{equation}
for every $w\in H^2(\Om)$ such that $\tr_0w=0.$
Here, $\tr_\nu w=\partial_\nu w|_{\partial\Om}$ is the standard normal-derivative trace of a $H^2$ function.

\newpage


\textbf{Well-defined, Unique, and Bounded}


Set
\[
    F(w):=\int_\Om(T_0u)\Delta w\dd x
          -\int_\Om\bigl(\Delta(T_0u)\bigr)w\dd x.
\]
Suppose $v_1,v_2\in H^2(\Om)$ satisfy
\[
    \tr_0v_1=\tr_0v_2=0,
    \qquad \tr_\nu v_1=\tr_\nu v_2.
\]
Then $w=v_1-v_2$ has both traces equal to zero. Under the usual smoothness
assumptions on $\partial\Om$,
\[
    H^2_0(\Om)
       =\{w\in H^2(\Om):\tr_0w=\tr_\nu w=0\},
\]
so $w\in H^2_0(\Om)$. For every $\varphi\in C_c^\infty(\Om)$, the
distributional Laplacian satisfies
\[
    \int_\Om(T_0u)\Delta\varphi\dd x
      =\int_\Om\bigl(\Delta(T_0u)\bigr)\varphi\dd x.
\]
Both sides are continuous in the $H^2$ norm, and
$H^2_0(\Om)=\overline{C_c^\infty(\Om)}^{\,H^2(\Om)}$.
Thus we can extend the identity to $w\in H^2_0(\Om)$, giving $F(w)=0$.
It follows that $F(v_1)=F(v_2)$, so the right-hand side of
\eqref{eq:T0weak} depends only on $\tr_\nu w$. 

On a smooth domain, the restricted normal trace
\[
    \tr_\nu:\{w\in H^2(\Om):\tr_0w=0\}
                       \longrightarrow H^{1/2}(\partial\Om)
\]
is surjective. Thus \eqref{eq:T0weak} defines a unique functional on
$H^{1/2}(\partial\Om)$. Furthermore, using Cauchy-Schwarz, $T_0u \in H^{-1/2}(\partial\Om).$ 

% \emph{Boundedness.}
% Cauchy--Schwarz gives
% \[
%     |F(w)|\leq C\bigl(\|T_0u\|_{L^2(\Om)}
%        +\|\Delta(T_0u)\|_{L^2(\Om)}\bigr)\|w\|_{H^2(\Om)}.
% \]
% Using a bounded lifting of the restricted normal trace, we obtain
% \[
%     |\dual{T_0u}{g}|
%       \leq C\|T_0u\|_{H(\Delta;\Om)}
%                  \|g\|_{H^{1/2}(\partial\Om)}.
% \]
% Hence
% \[
%     \boxed{T_0u\in H^{-1/2}(\partial\Om).}
% \]
% Strictly speaking, the notation on the left denotes the weak boundary
% trace of the interior quantity $T_0u$.

\subsection{The trace \texorpdfstring{$T_1u\in H^{-3/2}(\partial\Om)$}{T1}}

Now reconsider \eqref{eq:greens-second}, but impose
$\tr_\nu w=0$. This removes the boundary term containing $T_0u$ and
motivates
\begin{equation}\label{eq:T1weak}
    \dual{T_1u}{\tr_0w}_{H^{-3/2},H^{3/2}}
      :=\int_\Om\bigl(\Delta(T_0u)\bigr)w\dd x
        -\int_\Om(T_0u)\Delta w\dd x,
    % \quad
    % \begin{gathered}
    %    w\in H^2(\Om),\\[-0.2em]
    %    \tr_\nu w=0.
    % \end{gathered}
\end{equation}
for all $w\in H^2(\Om)$ such that $\tr_\nu w=0.$
The `well-defined' arguments are analogous to the $T_0u$ case. That is, the right hand side of \eqref{eq:T1weak} depends only on $w$ through $tr_0w.$ Furthermore,
\[
    \tr_0:\{w\in H^2(\Om):\tr_\nu w=0\}
                     \longrightarrow H^{3/2}(\partial\Om)
\]
is surjective. Thus, $T_1u$ is well-defined and unique. Finally, Cauchy-Schwarz tells us that $T_1u\in H^{-3/2}(\partial\Om).$
% \[
%     |\dual{T_1u}{g}|
%       \leq C\|T_0u\|_{H(\Delta;\Om)}
%                  \|g\|_{H^{3/2}(\partial\Om)}.
% \]
% Consequently,
% \[
%     \boxed{T_1u\in H^{-3/2}(\partial\Om).}
% \]
% Again, $T_1u$ denotes the weak boundary trace corresponding to the
% formal normal derivative $\partial_\nu(T_0u)$.

\newpage

\section{Some Conclusions / Generalisations}

The common approach is to use a Green's identity to isolate a not-so-well-defined boundary
term. In general, we start with 
\[
    a(u,v)=(Lu,v)_\Om+\dual{Bu}{\tr v},
\]
which motivates
\begin{equation}
    \dual{Bu}{\tr v}:=a(u,v)-(Lu,v)_\Om.
\end{equation}
Defining the trace then seems to follow three main steps:
\begin{enumerate}
    \item \textbf{Dependence on the boundary trace.}
          Show that the volume expression depends on $v$ only through
          $\tr v.$
    \item \textbf{Uniqueness.}
          Use surjectivity of the relevant test-function trace to
          obtain a unique boundary functional $Bu$.
    \item \textbf{Boundedness.}
          Apply Cauchy--Schwarz (and what else?) to show that
          $Bu$ is a dual element.
\end{enumerate}

For example, if $u\in H^1(\Om)$ and $\Delta u\in L^2(\Om)$, Green's
identity defines
\[
    \dual{\partial_\nu u}{\tr_0v}
     :=\int_\Om\grad u\cdot\grad v\dd x
        +\int_\Om(\Delta u)v\dd x.
\]
The distributional Laplacian establishes well-definedness; surjectivity
of $\tr_0$ establishes uniqueness; and Cauchy--Schwarz establishes $\partial_\nu u\in H^{-1/2}(\partial\Om)$.
Equivalently, because $\grad u\in H(\Div;\Om)$, we can use the
$H(\Div; \Om)$ normal trace of $\grad u$:
\[
    \partial_\nu u=\gamma_n(\grad u).
\]

\bibliography{Trace}
\bibliographystyle{plain}
\end{document}
